\documentclass[11pt,a4paper]{article}

% ---------------------------------------------------------------
% Packages
% ---------------------------------------------------------------
\usepackage[utf8]{inputenc}
\usepackage[T1]{fontenc}
\usepackage[margin=1in]{geometry}
\usepackage{hyperref}
\usepackage{url}
\usepackage{amsmath,amssymb}
\usepackage{listings}
\usepackage{xcolor}
\usepackage{booktabs}
\usepackage{array}
\usepackage[protrusion=true,expansion=false]{microtype}
\usepackage{parskip}
\usepackage{enumitem}
\usepackage{graphicx}
\usepackage{fancyhdr}
\usepackage{titlesec}
\usepackage[numbers,sort&compress]{natbib}

% ---------------------------------------------------------------
% Code listing style
% ---------------------------------------------------------------
\definecolor{codebg}{rgb}{0.97,0.97,0.97}
\definecolor{codeframe}{rgb}{0.85,0.85,0.85}
\definecolor{codecomment}{rgb}{0.4,0.4,0.6}
\definecolor{codestring}{rgb}{0.2,0.5,0.2}
\definecolor{codekw}{rgb}{0.1,0.1,0.7}

\lstset{
  backgroundcolor=\color{codebg},
  frame=single,
  rulecolor=\color{codeframe},
  basicstyle=\small\ttfamily,
  keywordstyle=\color{codekw}\bfseries,
  commentstyle=\color{codecomment}\itshape,
  stringstyle=\color{codestring},
  breaklines=true,
  breakatwhitespace=true,
  captionpos=b,
  numbers=none,
  xleftmargin=6pt,
  xrightmargin=6pt,
  aboveskip=8pt,
  belowskip=8pt,
  showstringspaces=false,
}

% ---------------------------------------------------------------
% Section formatting
% ---------------------------------------------------------------
\titleformat{\section}{\large\bfseries}{\thesection.}{0.5em}{}
\titleformat{\subsection}{\normalsize\bfseries}{\thesubsection.}{0.5em}{}
\titleformat{\subsubsection}{\normalsize\itshape\bfseries}{\thesubsubsection.}{0.5em}{}

% ---------------------------------------------------------------
% Header / footer
% ---------------------------------------------------------------
\pagestyle{fancy}
\fancyhf{}
\rhead{\small Subject-Enforced Completeness for AI Inspection Logs}
\lhead{\small Working paper, 2026}
\cfoot{\thepage}

% ---------------------------------------------------------------
% Hyperref config
% ---------------------------------------------------------------
\hypersetup{
  colorlinks=true,
  linkcolor=blue!60!black,
  citecolor=blue!60!black,
  urlcolor=blue!60!black,
  pdftitle={Subject-Enforced Completeness for AI Inspection Logs},
  pdfauthor={Atlas Project},
}

% ---------------------------------------------------------------
% Macros
% ---------------------------------------------------------------
\newcommand{\code}[1]{\texttt{#1}}
\newcommand{\term}[1]{\emph{#1}}
\newcommand{\todo}[1]{\textbf{\color{red}[TODO: #1]}}

% ---------------------------------------------------------------
\begin{document}
% ---------------------------------------------------------------

\title{\textbf{Subject-Enforced Completeness\\for AI Inspection Logs}}

\author{Atlas Project\\
\small \texttt{tomas.asin.gonzalez@gmail.com}\\[4pt]
\small Working paper --- \today\\
\small arXiv cs.CR (preprint)
}

\date{}
\maketitle

\begin{abstract}
\textit{Key transparency} systems such as CONIKS~\cite{coniks2015} and the IETF
\texttt{keytrans} working-group drafts~\cite{ietf-keytrans} allow users to monitor
their own entries in an append-only log and detect equivocation by the operator ---
without relying on global monitors. We transfer this pattern to a new domain:
\textbf{completeness of AI content-inspection logs}.

For a given subject's request stream, we show that completeness of the inspection log
--- whether every inspectable request was in fact inspected --- is \textbf{unilaterally
falsifiable} by the subject, without waiting for an external auditor. The mechanism
binds each inspection to a signed, monotonically-sequenced request emitted from the
subject's device. A gap in the subject's own sequence is a proof of omission: the
operator cannot skip an inspection and suppress the evidence unilaterally.

We implement the mechanism on RFC~9162~\cite{rfc9162} (Certificate Transparency v2),
extend it symmetrically to output inspection (every model response committed before
delivery, verified client-side), and provide a reference implementation. Layers~3--4
extend the scheme to external witnesses, split-view detection, and behavioral drift
observation. We state ten honest limits, including split-view exposure, geolocation
irresolvability, and covert post-inspection capability restrictions.

\textbf{We do not claim a new cryptographic primitive.} The contribution is
\textit{domain transfer}: carrying the subject-monitoring model of CONIKS and Key
Transparency from key-binding logs to AI content-inspection logs, and modelling
\textit{inspection-by-cause} as the invariant to enforce. Unlike prior work on
agent auditability --- which assumes honest infrastructure operators --- our mechanism
is designed specifically for the case where the operator controls both the AI pipeline
and the inspection log. The reference implementation comprises 1831 tests and is
mypy \texttt{-{}-strict} clean.
\end{abstract}

\tableofcontents
\newpage

% ---------------------------------------------------------------
\section{Introduction: Integrity $\neq$ Completeness}
% ---------------------------------------------------------------

Existing transparency infrastructure for AI systems focuses on \term{integrity}: proving
that log entries have not been altered after the fact. Merkle trees and signed tree heads
(STHs) provide this guarantee efficiently, and the Certificate Transparency ecosystem has
deployed it at scale.

Integrity, however, is the wrong property for inspection logs. An operator who wants to
hide that a request was processed without inspection does not alter existing log entries;
they simply \term{omit} the new one. A Merkle tree that grows correctly is consistent
with an operator who never inserts certain inspections. Verifying that the log has not
been tampered with says nothing about whether it is complete.

The gap matters structurally. In AI deployments, the entity that operates the model is
also the entity that operates the inspection infrastructure. This is not a malicious
arrangement --- it is the natural vertical integration of a model provider --- but it
creates a conflict of interest that integrity proofs do not address: the same party who
decides whether to run an inspection is the party who controls what the inspection log
records.

Our starting observation is that completeness, unlike integrity, has a per-subject
structure. A user who sends one hundred requests has a ground truth: one hundred signed
requests left their device. If the operator's log acknowledges only ninety-eight
inspections, two are unaccounted for. The user does not need a regulator, a second
operator, or a global monitor to compute this. They need their own counter.

The regulatory context sharpens the motivation. EU AI Act Articles~12, 13, and~53
impose record-keeping and transparency obligations on AI providers. When a regulator
asks ``was this request inspected?'', the current answer is ``trust the operator's
log.'' We ask whether the answer can be ``the subject's device can verify it
independently.''

\textbf{Fable~5 / next-generation frontier models} (Claude, GPT, Gemini) operate in a
regulatory grey zone: powerful enough to fall under systemic-risk obligations, deployed
widely enough that individual inspection gaps aggregate into policy-relevant failures.
This is deployment context motivating urgency, not a technical claim about any specific
model. We mention it in this section only.

% ---------------------------------------------------------------
\section{Threat Model}
% ---------------------------------------------------------------

\subsection{The Operator Is Automated Infrastructure}

A critical clarification: \textbf{the operator in this model is the AI provider's
automated infrastructure} (API gateway, inspection pipeline, log service), not an
individual human employee making per-request decisions. In a typical AI API deployment,
every request from every user flows through the same automated pipeline without manual
intervention.

This distinction matters for the threat model. Omissions arise from:

\begin{itemize}[noitemsep]
  \item \textbf{Configuration exceptions}: certain traffic patterns or API paths are
    routed around the inspection module.
  \item \textbf{Silent pipeline degradation}: the inspection service crashes or is
    rolled back; the gateway continues serving requests with a fallback that omits
    inspection.
  \item \textbf{Deliberate policy gaps}: the operator configures the system not to
    inspect certain categories of request by definition.
\end{itemize}

We model the operator as \textbf{semi-honest}: the infrastructure is operated by a
company that wants to demonstrate compliance (regulatory pressure, contractual
obligations), and the log infrastructure runs honestly for the requests it handles. The
adversarial capability is that the company controls which requests enter the inspection
pipeline at all.

The adversary's two goals:
\begin{enumerate}[noitemsep]
  \item \textbf{Selective omission}: process a request without creating an inspection
    record for it.
  \item \textbf{Split-view}: show the subject a log view that includes certain
    inspection records while showing a different view to the regulator.
\end{enumerate}

\subsection{What the Adversary Controls}

The operator controls: the inspection pipeline configuration; the log infrastructure;
and the client distribution channel.

\subsection{What the Adversary Does Not Control}

\begin{itemize}[noitemsep]
  \item \textbf{The subject's device-bound key.} A key generated inside the device's
    TPM or Secure Enclave (cf.\ Apple App Attest, Android Key Attestation) is
    hardware-bound and non-exportable.
  \item \textbf{The subject's local counter state.} The monotonic counter is maintained
    by the client on the subject's device.
  \item \textbf{The subject's local log replica.} The twin replica (\S\ref{sec:twin})
    is stored on the subject's device.
\end{itemize}

\subsection{What We Do Not Model}

We explicitly do not address: (a) an operator who is actively malicious and fully
non-compliant; (b) physical compromise of the subject's device; (c) a subject whose
``independent'' client was actually distributed by the operator (\S\ref{sec:limits},
limit~6.3); (d) content inspected out-of-band after retention without a new signed
request (limit~6.2).

% ---------------------------------------------------------------
\section{Mechanism: Subject-Enforced Completeness}
% ---------------------------------------------------------------

\subsection{Device-Bound Request Signing}

At first use, the subject authenticates via a federated identity provider (OAuth~2.0 /
OIDC). At registration time, the client silently generates an asymmetric key in the
device's hardware security module (TPM on PC/server, Secure Enclave on mobile). The
private key never leaves the hardware boundary.

Thereafter, every request the subject sends carries a \code{CosignedRequest}:

\begin{lstlisting}[language=Python]
CosignedRequest = {
    seq:          int,   # monotonic counter, starts at 0, never repeats
    payload_hash: str,   # SHA-256(prompt_bytes), hex
    signature:    str,   # device_key.sign( canonical_JSON({payload_hash, seq}) )
}
\end{lstlisting}

The canonical form signed is the compact, key-sorted JSON of
\code{\{payload\_hash, seq\}} --- no timestamps (timestamps are provided by the
operator's log entry, not by the subject, to avoid clock-skew disputes).

\textbf{Why not WebAuthn/Passkey per-request.} WebAuthn passkeys sign a server-provided
challenge in response to a user gesture; they cannot sign arbitrary payloads silently.
The correct primitive for silent per-request signing is a device-bound key in the
platform's HSM (Apple App Attest for mobile; client certificates bound to platform TPMs
for desktop).

\subsection{Inspection Binding and the In-Path Timing Guarantee}

For every \code{CosignedRequest} that triggers an inspection, the operator records:

\begin{lstlisting}[language=Python]
InspectionRecord = {
    seq:          int,   # from CosignedRequest
    payload_hash: str,   # SHA-256(payload) -- matches CosignedRequest.payload_hash
    cosig:        str,   # full CosignedRequest, JSON-serialised
    decision:     str,   # "allow" | "block" | "inspect"
    cause:        str,   # policy rule or heuristic that triggered this event
    timestamp_ns: int,   # operator-assigned epoch in nanoseconds
    salted_hash:  str,   # SHA-256(salt_i || payload) -- GDPR Art. 17 erasure hook
}
\end{lstlisting}

\textbf{Dual-hash leaf for GDPR Art.~17 (OSM-007).} The \code{InspectionRecord} carries
two hashes. \code{payload\_hash} is the hash the subject already signed; it enables
binding check~4 and is permanent. \code{salted\_hash} is derived from a 32-byte random
salt $s_i$ stored \emph{outside} the Merkle tree in a deletable store
(\code{SaltStore}). To exercise the right of erasure, the operator destroys $s_i$;
the committed salted hash becomes a dead-end without affecting tree integrity,
consistency proofs, or \code{detect\_omission()}.

\textbf{The critical invariant:} \textit{an inspection without an antecedent signed
request is invalid; a signed request without a corresponding inspection record is a
detectable omission.}

\textbf{In-path timing guarantee.} Because the filter is always-on and mandatory in the
API path, the \code{InspectionRecord} is committed to the Merkle log \emph{before} the
request is forwarded to the AI model. The model response cannot arrive at the client
before the commitment exists.

The API response carries proof back to the client:

\begin{lstlisting}[language=Python]
APIResponse = {
    result:            bytes,          # model response (opaque)
    seq_ack:           int,            # sequence the operator acknowledged
    leaf_bytes:        bytes,          # the committed InspectionRecord
    leaf_index:        int,
    inclusion_proof:   list[bytes],    # RFC 9162 inclusion proof
    sth:               SignedTreeHead, # {tree_size, root_hash, sig}
    consistency_proof: list[bytes],    # append-only proof from last head
    consistency_from:  int,
}
\end{lstlisting}

On each response the client runs six independent checks:

\begin{enumerate}[noitemsep]
  \item \textbf{STH signature} --- the head is authentically the operator's.
  \item \textbf{Append-only} --- the log was not rewritten since the client's last head.
  \item \textbf{Input inclusion} --- the committed \code{InspectionRecord} is in the tree.
  \item \textbf{Input binding} --- \code{payload\_hash} matches what the client signed.
  \item \textbf{Output inclusion} --- the committed \code{OutputInspectionRecord} is in the
    same tree.
  \item \textbf{Output binding} --- \code{SHA-256(result)} matches \code{output\_hash}
    in the committed output record.
\end{enumerate}

\subsection{\texttt{detect\_omission()}}

The subject maintains a local set of every \code{seq} value they have emitted:

\begin{lstlisting}[language=Python]
def detect_omission(observed_seqs: Sequence[int], last_emitted: int) -> list[int]:
    """Return seq values the operator did not acknowledge in [0, last_emitted]."""
    observed_set = set(observed_seqs)
    return [s for s in range(last_emitted + 1) if s not in observed_set]
\end{lstlisting}

This call requires no network request, no cooperation from the operator, and no
cryptographic computation beyond what the subject already holds locally. A non-empty
return value is a \term{unilateral proof of omission}.

\subsection{Twin Independent Log Replicas}
\label{sec:twin}

The operator maintains the authoritative Merkle log (RFC~9162). The subject maintains
a lightweight local replica: an ordered list of \code{(seq, operator\_root\_at\_seq)}
tuples, updated with each acknowledged inspection. Consistency proofs verify
append-only growth at each step.

\textbf{What this does not close.} The subject verifies their own view of the log. They
cannot prove that the regulator sees the same STH. Full closure requires external
witnesses (RFC~9162 \S5 gossip). The HTTP witness transport is implemented
(\code{HttpWitnessTransport} in \code{transparency/gossip.py}, Ed25519-verified
counter-signatures, \code{has\_quorum()}); deployment of independent witness nodes
remains infrastructure work outside this paper's scope. The residual split-view
exposure is stated as honest limit~\ref{limit:splitview}.

% ---------------------------------------------------------------
\section{Implementation}
% ---------------------------------------------------------------

The reference implementation lives in \code{src/atlas/transparency/} and is
intentionally decoupled from any specific AI product.

\textbf{\code{merkle\_tree.py}} --- RFC~9162 \S2.1 faithful. Domain-separation prefixes
prevent second-preimage attacks:
\begin{lstlisting}
leaf_hash(data)        = SHA-256(0x00 || data)
node_hash(left, right) = SHA-256(0x01 || left || right)
\end{lstlisting}
Exports \code{merkle\_root()}, \code{inclusion\_proof()} / \code{verify\_inclusion()},
\code{consistency\_proof()} / \code{verify\_consistency()}. All operations are pure
functions; the module has no I/O or state.

\textbf{\code{log.py}} --- \code{TransparencyLog} wraps the Merkle primitives with
append-only semantics, maintains the current STH, and signs it. Persistence is opt-in:
\code{TransparencyLog(path=Path(...))} writes each leaf as a base64-encoded line with
\code{fsync()} on append and reloads from disk on startup, preserving monotonic sequence
continuity across process restarts. A Read-API for deployers and regulators exposes the
log over HTTP: \code{GET /api/exec/api/v1/log/tree} (current STH),
\code{GET /api/exec/api/v1/log/entries} (leaf range, capped at 1000), and
\code{GET /api/exec/api/v1/log/proof/inclusion/\{i\}} (RFC~9162 inclusion proof),
providing the technical basis for EU AI Act Art.~26 deployer monitoring.

\textbf{\code{client\_cosign.py}} --- \code{ClientCosigner} emits monotonically-sequenced
\code{CosignedRequest} objects; \code{verify\_cosigned\_request()} validates signature
and payload hash; \code{detect\_omission()} implements \S3.3 above. The
\code{Signer} / \code{SigVerifier} interfaces accept any backend (software Ed25519 for
testing, hardware-bound keys in production).

\textbf{\code{witness.py}} --- \code{Witness.observe(sth)} accepts STHs from the
operator, verifies their signatures, and raises \code{SplitViewError} when two STHs with
the same \code{tree\_size} carry different \code{root\_hash} values.
\code{HttpWitnessTransport} submits STHs to independent witness endpoints over stdlib
\code{urllib}; counter-signatures are verified with Ed25519 before counting toward
quorum.

\textbf{Reference implementation and adversarial harness} (\code{docs/demo/},
reproducible, no network / no model). Ed25519 device-bound keys. Nine scenarios over a
5-request stream:

\begin{table}[h]
\centering
\small
\renewcommand{\arraystretch}{1.2}
\begin{tabular}{@{}lp{5.5cm}p{4cm}l@{}}
\toprule
\textbf{Session} & \textbf{Operator behaviour} & \textbf{Outcome} & \textbf{Ref.} \\
\midrule
A & Honest --- inspects + commits input and output & \code{detect\_omission() == []} & \S3 \\
B & Silently skips input inspection of seq=2 & \code{== [2]}, detected by subject alone & \S3.3 \\
C & Fakes an input ack (bogus STH + proof) & Rejected; \code{[2]} surfaces & \S3.2 \\
D & Rewrites a past log entry & Consistency proof fails & \S3.4 \\
E & Forges a request for a seq never sent & Rejected; signature $\neq$ registered key & \S2.3 \\
F & Signs receipt then omits; request lost in transit & Attributable omission isolated & \S6.8 \\
G & Commits input but omits output inspection seq=2 & Check 5 fails; cascade & \S3.2 \\
H & Shadow routing seq=2 (OSM-042) & Transparent to protocol & \S8.2 \\
J & Behavioral drift: rejection heuristics degrade & Canary delta flagged (probabilistic) & \S6.11 \\
\bottomrule
\end{tabular}
\caption{Adversarial harness scenarios. Sessions A--H plus J cover all paper claims.}
\end{table}

The harness is anchored in the test suite (\code{tests/test\_completeness\_demo.py},
\code{tests/test\_network\_reconciliation.py}, \code{tests/test\_output\_inspection.py});
1831 tests total, all passing.

% ---------------------------------------------------------------
\section{Related Work}
% ---------------------------------------------------------------

\textbf{Note on novelty.} The subject-monitoring mechanism is NOT novel. We cite its
predecessors as our foundation, not as work we improve upon.

\textbf{CONIKS}~\cite{coniks2015} is the direct ancestor. CONIKS allows users to monitor
their own key-binding entries in a transparency log, detecting inconsistency
\term{without global monitors} at low overhead ($<$20~kB/day). Client monitoring happens
automatically in the background, with no explicit user action required per check.
Over \term{key bindings}, CONIKS already delivers nearly everything we propose for
inspection logs: subject-side monitoring, automatic silent operation, and reduced
dependence on external witnesses.

\textbf{Key Transparency family.}~\cite{googkt2017,signalkt2023,appleckv2023,wakt2023}
Google Key Transparency (2017), Signal Key Transparency (2023), Apple Contact Key
Verification (2023), and WhatsApp Key Transparency (2023) are production deployments of
CONIKS-style subject-monitoring. The \textbf{IETF \texttt{keytrans} Working
Group}~\cite{ietf-keytrans} standardises this family. This is live production work
across major platforms; the primitive is mature and being standardised.

\textbf{Certificate Transparency.}~\cite{rfc6962,rfc9162} RFC~6962 and RFC~9162 provide
the append-only log infrastructure and gossip mechanisms for cross-operator STH
validation. Our implementation is RFC~9162 faithful. The split-view limit
(limit~\ref{limit:splitview}) is addressed by RFC~9162 \S5 gossip.

\textbf{The Attacker Moves Second}~\cite{attacker2024} demonstrates formally that
adaptive adversaries defeat $>$90\% of existing defences. This supports our decision not
to promise detection rates. We claim verifiability of the inspection record, not
detection probability.

\textbf{Our contribution, stated precisely.} The CONIKS/KT mechanism monitors
\term{key bindings}; the adversarial event is \term{equivocation}. We transfer this
model to \term{AI content-inspection logs}; the adversarial event is \term{omission}.
The invariant to protect changes: from ``my key binding is consistent across all views''
to ``every request I signed was preceded by a registered cause for inspection.'' We add
inspection-by-cause (\S3.2) as the specific record structure that makes this invariant
falsifiable. This is a domain-transfer and modelling contribution, not a cryptographic
primitive.

\subsection*{Concurrent Work (January--June 2026)}

\textbf{Notarized Agents}~\cite{sello2026} (Figuera, 2026) introduces receiver-side signing
for multi-agent workflows: each tool or service that receives an agent's call signs a Receipt
over what it actually observed, creating a Merkle-anchored audit trail. The threat model
is a \term{compromised agent} fabricating its own traces --- not a semi-honest operator
omitting inspections. The paper's honest-limits section explicitly acknowledges the
set-completeness gap it does not close: ``Owner receives verifiable individual receipts
but cannot cryptographically prove the log returned \textit{all} matching entries.
Operators could silently omit receipts while remaining cryptographically correct on
returned ones.'' This is exactly the gap our mechanism closes. The two contributions are
orthogonal in threat model: Sello secures agents against infrastructure they control; we
secure subjects against operators who control both the AI pipeline and the inspection log.

\textbf{Aegon}~\cite{aegon2026} (Baskaran et al., 2026) addresses a different domain:
\term{AI content licensing} (DRM). It extends JWT tokens with licensing claims, maintains
a CT-style Merkle ledger of licensing transactions, and adds Android StrongBox hardware
receipts for on-device compliance. The domain is licensing provenance, not safety
inspection; the threat model is content piracy, not operator omission. The paper is a
protocol design white paper; prototype implementation is stated as future work.
Despite the domain difference, the shared structure is notable: both systems use
append-only Merkle logs and hardware-bound receipts to make transactions undeniable.
The completeness gap --- guaranteeing that an adversarial SDK operator cannot silently
omit records --- is acknowledged as out of scope in Aegon; closing it is the central
contribution of our work.

\textbf{Auditable Agents}~\cite{auditableagents2026} (Nian et al., 2026) proposes a Merkle
log for agent execution traces, providing post-hoc auditability of tool calls and state
transitions. The mechanism assumes \term{honest tool providers}; completeness against a
dishonest operator is not addressed. Our contribution is precisely the case they assume away.

\textbf{Aegis}~\cite{aegis2026} (Mazzocchetti, 2026) combines ZK-STARK Proof-of-Conduct
with an ILK hash chain for governance of autonomous agents. Aegis addresses
\term{behavioral compliance} (did the agent follow rules?), not completeness of an
operator-side inspection log. The two mechanisms are complementary.

% ---------------------------------------------------------------
\section{Honest Limits}
\label{sec:limits}
% ---------------------------------------------------------------

These limits are not weaknesses to apologise for; they are the boundary conditions that
make every other claim credible.

\paragraph{6.1 Split-view: partial mitigation implemented, full closure pending.}
\label{limit:splitview}
\code{detect\_omission()} detects gaps in the subject's view of the log. It does not
prove that a regulator sees the same view. The HTTP witness transport is implemented
(\code{HttpWitnessTransport}, Ed25519-verified quorum, \code{has\_quorum()}); deployment
of independent witness nodes remains ecosystem work outside this paper's scope.

\paragraph{6.2 Out-of-band content and retained data.}
The mechanism covers the synchronous request path. Operator-side batch reprocessing of
stored prompts, asynchronous safety scans, or any inspection that does not reference a
live \code{CosignedRequest} are not covered.

\paragraph{6.3 Client-operator circularity.}
The guarantee rests on the independence of the client. If the operator both runs the AI
service and distributes the client software, they could ship a client that uses a weak
key or a counter that resets silently. The strong guarantee requires the client to be
distributed by a party independent of the AI operator.

\paragraph{6.4 Definition gaming.}
``What constitutes an inspection'' is defined by the operator. An operator can comply
with the mechanism while defining ``inspection'' narrowly enough to exclude most
meaningful safety checks. The mechanism provides a verifiable record of compliance with
the \term{defined} inspection policy; it does not enforce any particular policy.

\paragraph{6.5 Geolocation and export control are not resolvable in code.}
IP address, GPS, and round-trip latency are risk signals, not proof of geographic
location. What the filter makes verifiable is: ``we generated a risk score of X for
subject Y on date Z and escalated to KYC.'' Whether the KYC step resolves residency is
the legal system's responsibility.

\paragraph{6.6 Chain-of-thought is not auditable as ground truth.}
We log \code{(decision, action, cause)}, not the model's reasoning trace. Large language
model chain-of-thought outputs are known to be potentially post-hoc rationalisations. We
log the \term{act} and the \term{triggering rule}; the act is verifiable and the rule is
enumerable. The reasoning is not.

\paragraph{6.7 Retroactive compliance: architectural mitigation, not cryptographic closure.}
The in-path guarantee is a property of the deployed architecture, not a cryptographic
invariant. An operator who deliberately engineers a bypass --- letting the model run
before the filter commits, then inserting a retroactive record --- can defeat it without
forging any signature.

\paragraph{6.8 Network failures create plausible deniability for unconfirmed requests.}
A gap in the subject's sequence has two indistinguishable causes: operator omission, or
the request never arrived. The mitigation (OSM-040) is a signed receipt: the operator
returns a \code{Receipt\{seq, payload\_hash, received\_at\_ns\}} signed on arrival,
\emph{before} inspection. A subject who holds a valid receipt for \code{seq=n} but
finds no inclusion proof for \code{n} has an \term{attributable omission}. Exercised in
Session~F and \code{tests/test\_network\_reconciliation.py}.

\paragraph{6.9 Legal shield is a hypothesis without legal counsel.}
Section~\ref{sec:deployment} proposes that verifiable inspection logs create a mutual
protection incentive. Whether this argument succeeds in any jurisdiction depends on local
law. We present it as a plausible deployment incentive, not a legal fact.

\paragraph{6.10 Covert post-inspection capability restrictions are not detectable.}
The completeness mechanism ensures every request was inspected; it cannot constrain what
the model does \emph{after} inspection returns ``allow.'' On 2026-06-10, Anthropic
launched Claude Fable~5 with an undisclosed instruction that silently reduced response
quality for certain AI-research queries without informing the user; the instruction was
reversed within 24~hours after public disclosure. The episode confirms that a complete
and verifiable inspection log --- one that recorded ``allow'' for each of those requests
--- would not have surfaced the degradation. Our protocol makes completeness of
inspection verifiable; it makes no claim about the faithfulness of the model's subsequent
output.

\paragraph{6.11 Behavioral Drift Detection (OSM-054) --- Limitations and Open Problems.}
\label{limit:drift}
Three angles are proposed: (A)~behavioral delta observation via canary prompts;
(B)~ex-post consistency proofs; (C)~shadow divergence. All three face common barriers:
capability masking can be implicit in model weights; false positives are likely when the
model improves legitimately; false negatives occur when restrictions operate at
sub-threshold probability levels. We state behavioral drift detection as a
\textbf{future research problem} and an \textbf{open investigative tool}, not a closed
detection method. See \code{src/atlas/security/behavioral.py} for the reference
implementation.

% ---------------------------------------------------------------
\section{Deployment and Incentives}
\label{sec:deployment}
% ---------------------------------------------------------------

\subsection{Mutual Protection}

The verifiable inspection log creates a protection asymmetry that benefits both parties.
The operator can produce an inclusion proof when a subject claims harm from AI-generated
content: ``inspection record for \code{(seq=n, payload\_hash=h)} is in the log at
position~$k$, with STH signed at time~$t$.'' The subject can audit the \code{cause}
field when they claim wrongful blocking. The same log that protects the operator from
``you never inspected this'' also exposes the operator to ``you inspected this without
cause.'' Both parties benefit if and only if the operator is operating honestly within
their declared policy.

\subsection{EU AI Act Alignment}

\begin{table}[h]
\centering
\small
\renewcommand{\arraystretch}{1.2}
\begin{tabular}{@{}lp{8cm}@{}}
\toprule
\textbf{Article} & \textbf{Mechanism contribution} \\
\midrule
Art.~9 (risk management) & Cause-based inspection (\S3.2) creates auditable per-request risk assessment record. \\
Art.~12 (record-keeping) & Merkle log + \code{InspectionRecord} (decision + cause + timestamp); append-only, tamper-resistant, unlimited retention. \\
Art.~13 (transparency) & Subject-held replica + \code{detect\_omission()} enable individual-request-level transparency without trusting operator. \\
Art.~14 (human oversight) & Verifiable log makes inspection record independently auditable by regulator or subject. \\
Art.~26 (deployer obligations) & Read-API (\code{GET /api/exec/api/v1/log/\ldots}) exposes per-request evidence base for deployer monitoring. \\
Art.~53 (systemic risk) & Per-request inspection evidence at scale supports documentation obligation. \\
\bottomrule
\end{tabular}
\caption{EU AI Act alignment. The filter provides technical evidence; it does not substitute for substantive legal obligations.}
\end{table}

\subsection{Enforcement Boundary}

Breaking the chain --- refusing to participate in the protocol --- disables the guarantee
for that session. The preferred design is \term{graceful degradation}: chain breaks are
logged and flagged for review, not used as hard gates, to avoid converting every network
anomaly into a service outage.

% ---------------------------------------------------------------
\section{Conclusion}
% ---------------------------------------------------------------

We have shown that completeness of AI content-inspection logs is \textit{unilaterally
falsifiable} by the subject: a gap in the subject's monotonically-signed request sequence
is a proof of omission that requires no external auditor and no cooperation from the
operator. The mechanism transfers the subject-monitoring model of CONIKS and Key
Transparency to a new domain, adds inspection-by-cause as the specific invariant to
protect, and extends it symmetrically to output inspection. Concurrent work on agent
auditability (Notarized Agents, Auditable Agents, Aegon, Aegis) addresses related but
orthogonal problems; none closes the set-completeness gap against a semi-honest operator
that this mechanism closes.

\textbf{Open work.} Three threads remain outside this paper's scope.
(1)~\textit{Independent witness deployment}: the HTTP witness transport is implemented;
operating independent witness nodes requires ecosystem coordination beyond a single
deployment.
(2)~\textit{Behavioral faithfulness}: proving that a model's observable behavior on
canary inputs predicts its behavior on the full production distribution requires
interpretability infrastructure not yet available at API level.
(3)~\textit{KYC / export-control binding}: the mechanism provides a device-bound
signing key; resolving legal residency for export-controlled access requires a KYC
step that is operational and legal, not cryptographic.

% ---------------------------------------------------------------
\appendix
\section{Four-Layer Defense Stack}
\label{sec:defensestack}
% ---------------------------------------------------------------

The completeness mechanism (§§3--6) does not exist in isolation. This appendix
synthesises the four-layer stack, states what each layer does not guarantee, and
identifies the open research problems no layer currently closes.

The completeness mechanism is one layer of a coordinated defense-in-depth stack.

\subsection{Layer 1--2: Completeness Verification (Input + Output)}

Enforces symmetric completeness on both the input request and the output response.
Omissions are detectable by the subject through \code{detect\_omission()} without
relying on the operator. The structural invariants are OSM-007 (input completeness) and
OSM-042 (output completeness). Both are closed.

\textbf{Does not guarantee:} fabricated cause fields; split-view attacks.

\subsection{Layer 3: Active Defense and Consensus (OSM-031/052/053)}

External-witness network following RFC~9162 gossip. A campaign metric
(\code{C\_attempts}, \code{K\_attribution}: $\geq 3$ consecutive attempts with cosine
similarity $>0.7$) triggers flagging with co-signed identity trail.

\textbf{Does not guarantee:} per-attempt detection; single-operator witness deployment.

\subsection{Layer 4: Behavioral Drift Observation (OSM-054)}

A canary framework: fixed probes re-issued periodically; drift in response distribution
signals undisclosed model changes. Session~J (demo) exercises this. Probabilistic; does
not attribute drift to adversarial manipulation vs.\ legitimate update.

\subsection{Synthesis}

\begin{table}[h]
\centering
\small
\renewcommand{\arraystretch}{1.2}
\begin{tabular}{@{}lp{5cm}p{5cm}@{}}
\toprule
\textbf{Layer} & \textbf{Threat closed} & \textbf{Threat not closed} \\
\midrule
L1--L2 (OSM-007/042) & Inspection omission detectable by subject & Fabricated cause; split-view \\
L3 (OSM-031/052/053) & Split-view via witnesses; campaign attribution & Per-attempt detection; single-operator witness \\
L4 (OSM-054) & Behavioral drift early warning (canary, heuristic) & Behavioral faithfulness; covert capability detection \\
\bottomrule
\end{tabular}
\caption{Orthogonal defense layers. No layer subsumes another.}
\end{table}

\textbf{Open research.} Two problems remain outside the scope of this paper and the
published literature as of June~2026: (1)~\term{behavioral faithfulness} --- proving
that a model's behavior on a finite canary set predicts its behavior on the full
production input distribution; (2)~\term{covert capability detection} --- detecting that
a model has acquired a capability without access to its weights or training process.

% ---------------------------------------------------------------
\section*{Acknowledgements}
% ---------------------------------------------------------------

Reference implementation: \url{https://github.com/therealronin23/atlas} (private;
available upon request for reviewer verification). All claims are anchored in 1831
automated tests, mypy \texttt{-{}-strict} clean.

% ---------------------------------------------------------------
% Bibliography
% ---------------------------------------------------------------
\bibliographystyle{plainnat}
\bibliography{paper_subject_enforced_completeness}

\end{document}
